\chapter*{Abstract}

UMBRAL is a behavioral monitoring system for Argentine small and medium-sized businesses that addresses the post-authentication blind spot: an attacker with valid credentials operates with the legitimate user's permissions without triggering rule-based controls. Following the separation between policy decision point and policy enforcement point in NIST SP 800-207, a lightweight Windows agent collects process, file and connection metadata and sends it over TLS to a central API. The API computes a risk score from 0 to 100 for each event, with continuity across sessions and temporal decay, and returns a graduated action that the agent executes on the device: from an alert to terminating processes or reversibly denying access to folders. A web dashboard shows each user's risk with its explanation and lets administrators adjust thresholds and actions at run time.

The model combines Isolation Forest over daily activity with an explicit location rule and resource novelty detection, trained on real logs from the CLUE-LDS dataset. On real days, the rate of days at high level is 1.81\,\% on the development set and 1.05\,\% in a single evaluation on the held-out set, and the 95th percentile of computation latency is 13.5~ms. The 801 automated tests exceed 85\,\% line coverage. Detection capability could not be measured, because no dataset in the domain includes real labeled attacks, and the model trained on cloud data flags 53.66\,\% of the days of a Windows host domain, compared with 4.13\,\% for a model trained in that domain. Collecting the system's own telemetry and obtaining labeled incidents are therefore the main future work. The economic analysis shows viability in a moderate adoption scenario.

\vspace{1em}
\noindent\textbf{Keywords:} UEBA, Zero Trust, anomaly detection, Isolation Forest, risk scoring, endpoint response, SMB.
